\documentclass[../veristore_security_model.tex]{subfiles}

\begin{document}
\section{External Security Notes}

% COMMENT: Document any security-relevant information that should be communicated to users,
% deployers, or other teams. This includes known limitations, security best practices for
% deployment, configuration recommendations, and any threats that are out of scope for veri-store
% but users should be aware of. Examples: recommend using TLS for network communication,
% suggest server isolation strategies, note that data confidentiality requires client-side encryption.

\subsection{Deployment Recommendations}

We recommend a simple but safe deployment setup:
\begin{itemize}
    \item Keep the five storage servers on a private network and do not expose them directly to the public internet.
    \item Use HTTPS through a reverse proxy if clients are not running on localhost.
    \item Set a strong shared API token and rotate it if it is leaked.
    \item Run each server process with limited filesystem permissions so it can only access its own data folder.
    \item Turn on request logging so failed verification events are easy to trace during debugging and demos.
\end{itemize}

\subsection{Known Limitations}

% COMMENT: Document security limitations that users should be aware of. What attacks does the
% system NOT protect against? What trust assumptions are made? Under what conditions might
% security guarantees not hold?

\subsection{Integration Considerations}

% COMMENT: If veri-store is used as a component in a larger system, what security properties
% does it provide and what must be handled by the integrating system? For example: veri-store
% ensures integrity but not confidentiality; authentication must be provided by the application.

\newpage
\end{document}
